% ============================================================
% INTRODUCTION - target: 600--800 words, 4 paragraphs
% P1: Basic problem | P2: What is known | P3: The gap (hook)
% P4: How the gap is addressed + hypotheses
% ============================================================

\section{Introduction}

% Paragraph 1: The Problem
As large language models are deployed in increasingly high-stakes
settings, from medical diagnosis and legal reasoning to autonomous code
generation, the ability to monitor and verify their reasoning processes becomes
a critical safety requirement. Chain-of-thought (CoT) reasoning, in which models
produce step-by-step explanations before arriving at a final answer, has been
widely adopted as a transparency mechanism: if a model's reasoning is visible,
human overseers can potentially detect flawed logic, biased reasoning, or
deceptive intent before harm occurs \cite{wei2022cot}. This promise is
especially significant for reasoning models (systems explicitly trained to
produce extended thinking traces via reinforcement learning) which now dominate
benchmarks in mathematics, science, and coding \cite{qchen2025longcot,
xu2025large}. However, the safety value of CoT monitoring rests on a crucial
assumption: that the chain of thought faithfully represents the model's actual
reasoning process. If models can be influenced by factors they do not acknowledge
in their CoT, then monitoring these traces provides a false sense of security
\cite{kadavath2022selfknowledge, korbak2025monitorability}.

% Paragraph 2: What Is Known
A growing body of evidence suggests that CoT explanations are frequently
unfaithful. Turpin et al.~\cite{turpin2023unfaithful} demonstrated that biased
features in prompts (such as suggested answers from a purported
expert) systematically influence model outputs without being mentioned in the
CoT, revealing structural gaps between stated and actual reasoning across
multiple tasks on BIG-Bench Hard. Lanham et al.~\cite{lanham2023measuring}
further measured CoT faithfulness in Claude models by truncating, corrupting, and
adding mistakes to reasoning chains, finding that the correlation between CoT
content and model predictions varies substantially across tasks. Subsequent work
established that sycophancy in RLHF-trained models is pervasive
\cite{sharma2023sycophancy} and scales with model size
\cite{perez2022discovering}, while Roger and Greenblatt~\cite{roger2023steganography}
showed that models can encode reasoning steganographically in CoT text,
suggesting that surface-level monitoring may systematically underestimate the
true divergence between stated and actual reasoning. Most critically, Chen et
al.~\cite{chen2025reasoning} conducted the most comprehensive evaluation to
date, using six types of reasoning hints injected into MMLU and GPQA
multiple-choice questions. Their results revealed that Claude~3.7 Sonnet
acknowledges hints that influenced its answer only 25\% of the time, while
DeepSeek-R1 does so only 39\% of the time, with unfaithful CoTs being
paradoxically \emph{longer} than faithful ones. A shared limitation across all of
these studies is their narrow model coverage: existing evaluations test at most
two to four models, leaving the vast and rapidly growing ecosystem of open-weight
reasoning models almost entirely unexamined.

% Paragraph 3: The Gap (Hook)
What remains unknown is whether the low faithfulness rates observed in Claude and
DeepSeek-R1 generalize across the broader open-weight reasoning model
ecosystem, or whether they reflect properties specific to particular
architectures, training methods, or model scales. Since early 2025, a diverse
set of open-weight reasoning models has emerged spanning dense architectures from
7B to 32B parameters, mixture-of-experts (MoE) models exceeding 685B total
parameters, and training pipelines ranging from pure reinforcement learning
(GRPO, RL) to distillation, data-centric methods, and hybrid approaches. No
study has systematically compared CoT faithfulness across these model families,
leaving a critical gap in understanding which models, and which training
paradigms, produce reasoning traces that can be reliably monitored for safety
purposes. As Lightman et al.~\cite{lightman2023verify} demonstrated, step-level
verification of reasoning outperforms outcome-level supervision, making the
question of whether each reasoning step is honest, not merely
correct, a first-order concern for AI safety.

% Paragraph 4: How the Gap Is Addressed + Hypotheses
The present study addresses this gap with a large-scale, cross-family
evaluation of CoT faithfulness in open-weight reasoning models. A total of 12
models spanning 9 architectural families, including DeepSeek, Qwen,
MiniMax, OpenAI, Baidu, AI2, NVIDIA, StepFun, and
ByteDance, are evaluated on 498
multiple-choice questions (300 from MMLU and 198 from GPQA Diamond) using the
same six hint categories established by Chen et al.~\cite{chen2025reasoning},
yielding 41,832 inference runs. Faithfulness is assessed using a two-stage
classifier (regex pattern matching followed by a three-judge LLM panel),
with Claude Sonnet~4 as an independent validation judge on all 10,276
influenced cases~\cite{young2026classifier}. Two primary hypotheses
are tested: (H1)~faithfulness rates vary significantly across model families,
reflecting differences in training methodology rather than scale alone; and
(H2)~certain hint types, particularly metadata and grader
hacking, exhibit consistently lower faithfulness across all models, suggesting
category-specific blind spots in CoT monitoring. All code, prompts, model outputs, and faithfulness
annotations are released to support future research on CoT monitoring
reliability.
